% chktex-file 44 32 36 38 18 11 8 26

\pagestyle{DS_FP}
%\NomPrenomNote{}

\bigskip

\exods{0}
%%% POLYNESIE 2025 - J1 - exo1

\section*{PARTIE A}

\begin{enumerate}[series=monDS]
    \item Donner l'ordre de parcours des sommets du graphe.\hfill \reponseLigne[3cm]{ABCEDF}

    \item Donner la table de routage de F à l’initialisation du protocole RIP.

    \begin{blocvisible}[height=1.5]

\begin{center}
\begin{tabular}{|c|c|c|}
\hline
\textbf{Routeur} & \textbf{Nombre de sauts} & \textbf{Prochain routeur} \\
\hline
\textbf{D} & 1 & -- \\
\hline
\end{tabular}
\end{center}
\end{blocvisible}


    \item Donner la table de routage de A après une première itération de RIP (deux réponses sont possibles).
    
     \begin{blocvisible}[height=3.]

\begin{center}
\begin{tabular}{|c|c|c|}
\hline
\textbf{Routeur} & \textbf{Nombre de sauts} & \textbf{Prochain routeur} \\
\hline
\textbf{B} & 1 & -- \\
\hline
\textbf{C} & 1 & -- \\
\hline
\textbf{E} & 1 & -- \\
\hline
\textbf{D} & 2 & E ou C \\
\hline
\end{tabular}
\end{center}
\end{blocvisible}
    \item Donner le numéro de l’itération de RIP à partir duquel les tables des routeurs du réseau ne varient plus. \hfill 
\begin{blocvisible}[height=1]
    Les tables de routage ne varient plus à partir de la troisième itération de RIP.

\end{blocvisible}



    \item Donner la nouvelle table de routage de A après stabilisation de RIP (deux réponses sont possibles).

    \begin{blocvisible}[height=3.5]
\begin{center}
\begin{tabular}{|c|c|c|}
\hline
\textbf{Routeur} & \textbf{Nombre de sauts} & \textbf{Prochain routeur} \\
\hline
\textbf{B} & 1 & -- \\
\hline
\textbf{C} & 1 & -- \\
\hline
\textbf{E} & 1 & -- \\
\hline
\textbf{D} & 2 & E ou C \\
\hline
\textbf{F} & 3 & E ou C\\
\hline
\end{tabular}
\end{center}
\end{blocvisible}
\end{enumerate}


\section*{PARTIE B}

    \begin{enumerate}[series=monDS, resume]
        \item Compléter le code de la méthode \verb|relie|.

        \begin{blocvisible}[height=4.5]

    \begin{CodePythontexAlt}[TaillePolice=\large, Largeur=1\linewidth,Centre,Lignes=true,PremLigne=1]{}
def relie(self, autre): 
    if autre not in self.voisins: 
        self.voisins.append(autre) 
        self.nb_sauts[autre] = 1 
        self.prochain[autre] = None 
        if self not in autre.voisins: 
            autre.relie(self)
    \end{CodePythontexAlt}
\end{blocvisible}


        \item Écrire la méthode \verb|relie_liste| de la classe \verb|Routeur| qui prend en paramètre une liste de routeurs \verb|lst| et qui relie le routeur \verb|self| à chacun des routeurs de la liste \verb|lst|.
        
        \begin{blocvisible}[height=2.5]
    \begin{CodePythontexAlt}[TaillePolice=\large, Largeur=1\linewidth,Centre,Lignes=true,PremLigne=1]{}
def relie_liste(self, lst):
    for autre in lst: 
        self.relie(autre)
    \end{CodePythontexAlt}
\end{blocvisible}


    
\item Écrire les instructions manquantes pour relier les routeurs de manière à obtenir le graphe de la figure \ref{reseau_exo1}.

\begin{blocvisible}[height=2.5]
    \begin{CodePythontexAlt}[TaillePolice=\large, Largeur=1\linewidth,Centre,Lignes=true,PremLigne=1]{}
A.relie_liste([B, C, E]) 
D.relie_liste([C, E, F])
C.relie(E) 

\end{CodePythontexAlt}

\end{blocvisible}

    \item Compléter le code de la méthode \verb|met_a_jour_table| ci-dessus.
    
    \begin{blocvisible}[height=4.5]
    \begin{CodePythontexAlt}[TaillePolice=\large, Largeur=1\linewidth,Centre,Lignes=true,PremLigne=1]{}
def met_a_jour_table(self, autre):
    for r in autre.nb_sauts: 
        if r != self: 
            if (r not in self.nb_sauts or self.nb_sauts[r] > \
                                        autre.nb_sauts[r] + 1): 
                self.nb_sauts[r] = autre.nb_sauts[r] + 1 
                self.prochain[r] = autre.nom
\end{CodePythontexAlt}
\end{blocvisible}

\medskip
    \item Écrire la méthode \verb|itere_rip| qui prend en paramètre le routeur \verb|self| et met à jour sa table de routage lorsqu’il reçoit la table de routage de chacun des routeurs présents dans la liste de ses voisins.
\begin{blocvisible}[height=2.5]
    \begin{CodePythontexAlt}[TaillePolice=\large, Largeur=1\linewidth,Centre,Lignes=true,PremLigne=1]{} 
def itere_rip(self): 
    for autre in self.voisins: 
        self.met_a_jour_table(autre)
\end{CodePythontexAlt}
    
\end{blocvisible}
    \medskip
    \item Écrire une fonction qui prend en paramètre une liste de routeurs \verb|l_routeurs| et qui réalise une itération du protocole RIP pour tous les routeurs de \verb|l_routeurs|.
\begin{blocvisible}[height=2.5]
    \begin{CodePythontexAlt}[TaillePolice=\large, Largeur=1\linewidth,Centre,Lignes=true,PremLigne=1]{} 
def iteration_rip(l_routeurs):
    for r in l_routeurs:
        r.itere_rip()
\end{CodePythontexAlt}
    
\end{blocvisible}

    \item Écrire une version modifiée du code de la méthode \verb|itere_rip| de la classe \verb|Routeur| de telle sorte que celle-ci renvoie \verb|True| dans le cas où le routeur \verb|self| a procédé à une modification de sa table de routage au cours de lexécution de la méthode \verb|itere_rip|, et \verb|False| sinon.

\begin{blocvisible}[height=4.]
    \begin{CodePythontexAlt}[TaillePolice=\large, Largeur=1\linewidth,Centre,Lignes=true,PremLigne=1]{}
def itere_rip(self): 
    modifie = False
    for autre in self.voisins: 
        if self.met_a_jour_table(autre):
            modifie = True
    return modifie
\end{CodePythontexAlt}
\end{blocvisible}

\begin{minipage}{0.44\textwidth}
        \item Compléter le code du programme principal afin que celui-ci mette à jour les tables de routage des routeurs présents dans la liste \verb|liste_routeurs| jusqu’à ce qu’il ne soit plus nécessaire de faire des mises à jour des tables de routage.

\end{minipage}\hfill
\begin{minipage}{0.48\textwidth}


\begin{blocvisible}[height=4.]
    \begin{CodePythontexAlt}[TaillePolice=\large, Largeur=1\linewidth,Centre,Lignes=true,PremLigne=1]{}

modifie = True
while modifie:
    modifie = False
    for r in liste_routeurs:
        if r.itere_rip():
            modifie = True
\end{CodePythontexAlt}
\end{blocvisible}
\end{minipage}
\end{enumerate}